\chapter{Rupert pairs}
\label{chap:prelims}

Write $\pi \colon \R^3 \to \R^2$ for the projection $(x_0, x_1, x_2) \mapsto (x_0, x_1)$
onto the screen.

\begin{definition}
\label{def:rupert}
\lean{IsRupert}
\leanok
A convex polyhedron with vertex set $V$ is \emph{Rupert} if there are rotations
$R_i, R_o \in SO(3)$ and an offset $o \in \R^2$ such that
\[
  o + \pi(R_i \hull{V}) \subseteq \operatorname{int} \pi(R_o \hull{V}).
\]
\end{definition}

\begin{definition}
\label{def:exists_non_rupert}
\lean{ExistsNonRupertPolyhedron}
\leanok
\uses{def:rupert}
The main statement: there is a finite set $V \subseteq \R^3$ whose convex hull has
nonempty interior and which is not Rupert.
\end{definition}

\section{Rupert sets}

\begin{theorem}[Rupert Polyhedron iff Rupert Set]
\label{thm:rupert_iff_rupert_set}
\lean{rupert_iff_rupert_set}
\leanok
\uses{def:rupert}
The convex polyhedron with vertex set $V$ is Rupert if and only if the convex hull of
$V$ is a Rupert set, i.e.\ there is a pose (two rotations and an offset) under which
the closure of the inner shadow lies in the interior of the outer shadow.
\end{theorem}
\begin{proof}
\leanok
The convex hull of a finite set is compact, so its shadow is closed.
\end{proof}

\begin{theorem}
\label{thm:rupert_implies_rot_rupert}
\lean{rupert_implies_rot_rupert}
\leanok
\uses{def:pointsymmetric}
Let $S \subseteq \R^3$ be point-symmetric and convex. If a pose with offset $o$
passes $S$ through itself, then so does the same pose with offset $0$.
\end{theorem}
\begin{proof}
\leanok
Write $K$ for the closure of the inner shadow and $L$ for the interior of the outer
shadow; both are point-symmetric and $L$ is convex. From $o + K \subseteq L$ and
symmetry, $-o + K = -(o + K) \subseteq L$. Every $k \in K$ is the midpoint of
$o + k \in L$ and $-o + k \in L$, so $k \in L$ by convexity.
\end{proof}

\section{Shadows and the separation criterion}

From now on $\PPP$ denotes the Noperthedron (Definition~\ref{def:noperthedron}).

\begin{definition}
\label{def:shadow}
\lean{Cayley.shadow, Cayley.RupPair}
\leanok
\uses{def:noperthedron}
The \emph{shadow} of a $3 \times 3$ matrix $R$ is $\mathrm{Sh}(R) = \pi(R\,\PPP)$.
A pair of matrices $(R_o, R_i)$ is a \emph{Rupert pair} if
$\overline{\mathrm{Sh}(R_i)} \subseteq \operatorname{int} \mathrm{Sh}(R_o)$.
\end{definition}

\begin{lemma}
\label{lem:exists_rup_pair}
\lean{Cayley.exists_rupPair}
\leanok
\uses{def:shadow, thm:rupert_iff_rupert_set, thm:rupert_implies_rot_rupert, lemma:nopert_point_symmetric}
If $\PPP$ is Rupert, then there is a Rupert pair $(R_o, R_i)$ with
$R_o, R_i \in SO(3)$.
\end{lemma}
\begin{proof}
\leanok
Combine Theorem~\ref{thm:rupert_iff_rupert_set} and
Theorem~\ref{thm:rupert_implies_rot_rupert}; $\PPP$ is point-symmetric
(Lemma~\ref{lemma:nopert_point_symmetric}) and convex.
\end{proof}

For $w = (w_0, w_1) \in \R^2$ write $\langle w, x \rangle = w_0 x_0 + w_1 x_1$ for
$x \in \R^3$ (only the screen coordinates enter).

\begin{theorem}[Separation criterion]
\label{thm:separation}
\lean{Cayley.not_rupPair_of_sep}
\leanok
\uses{def:shadow}
Let $S$ be a vertex of $\PPP$ and $w \ne 0$. If
\[
  \langle w, R_o T \rangle \le \langle w, R_i S \rangle
\]
for every vertex $T$ of $\PPP$, then $(R_o, R_i)$ is not a Rupert pair.
\end{theorem}
\begin{proof}
\leanok
The outer shadow is the convex hull of the points $\pi(R_o T)$, so it lies in the
closed half-plane $\{ z : \langle w, z \rangle \le \langle w, R_i S \rangle \}$. The
point $\pi(R_i S)$ of the inner shadow lies on the boundary line of this half-plane,
so it is not an interior point of the outer shadow.
\end{proof}

Only a non-strict inequality is needed, so poses in which the two shadows touch need
no special treatment.
